% ============================================================
%  CHAPTER 3: KINEMATICS
%  The Physics Codex: A Complete University Foundation
%  Korede Samuel Omotosho (Falcon)
%  Aurikrex Learning Systems, 2026
%
%  File: chapters/part1/ch03_kinematics.tex
% ============================================================

\chapter{Kinematics}
\label{ch:kinematics}

\chapterquote{Nature uses only the longest threads to weave
her patterns, so each small piece of her fabric reveals
the organisation of the entire tapestry.}{Richard Feynman}

\begin{objectivebox}
By the end of this chapter, you should be able to:
\begin{enumerate}[noitemsep]
  \item Define and distinguish between distance,
  displacement, speed, velocity and acceleration
  \item Derive and apply the four equations of uniform
  acceleration
  \item Analyse motion using displacement-time and
  velocity-time graphs
  \item Solve problems involving bodies falling freely
  under gravity
  \item Analyse projectile motion by resolving into
  horizontal and vertical components
  \item Calculate range, time of flight and maximum
  height of projectiles
  \item Describe and solve problems in relative motion
\end{enumerate}
\end{objectivebox}

\vspace{0.4cm}

\minitoc

\vspace{0.5cm}

% ============================================================
\section{What is Kinematics?}
% ============================================================

\textbf{Kinematics} is the branch of mechanics that
describes \textit{how} objects move --- their position,
speed, direction and change of speed --- without asking
\textit{why} they move. The \textit{why} is the business
of dynamics (Chapter 4). Kinematics is the language of
motion; dynamics provides the explanation.

To understand kinematics properly, we need precise
definitions of the terms we use. In everyday speech,
``distance'' and ``displacement'' mean much the same
thing. In Physics they are completely different quantities.

% ============================================================
\section{Fundamental Definitions}
% ============================================================

\begin{definitionbox}[Distance and Displacement]
\textbf{Distance} is the total length of the path
travelled by a body, regardless of direction.
It is a \textbf{scalar} quantity. SI unit: metre (m).

\vspace{0.4cm}

\textbf{Displacement} is the straight-line distance from
the initial position to the final position, measured in
a specified direction. It is a \textbf{vector} quantity.
SI unit: metre (m).
\end{definitionbox}

\vspace{2.3cm}

\begin{definitionbox}[Speed and Velocity]
\textbf{Speed} is the rate of change of distance with
time. It is a \textbf{scalar}. SI unit: m\,s$^{-1}$.
\[
  \text{Speed} = \frac{\text{distance}}{\text{time}}
\]

\textbf{Velocity} is the rate of change of displacement
with time. It is a \textbf{vector}. SI unit: m\,s$^{-1}$.
\[
  \vect{v} = \frac{\Delta\vect{s}}{\Delta t}
\]

\textbf{Average velocity} over a time interval $\Delta t$:
\[
  \bar{v} = \frac{\text{total displacement}}
  {\text{total time}} = \frac{\Delta s}{\Delta t}
\]
\end{definitionbox}

\begin{definitionbox}[Acceleration]
\textbf{Acceleration} is the rate of change of velocity
with time. It is a \textbf{vector}. SI unit: m\,s$^{-2}$.
\[
  \vect{a} = \frac{\Delta\vect{v}}{\Delta t}
  = \frac{\vect{v} - \vect{u}}{t}
\]
where $\vect{u}$ is the initial velocity and $\vect{v}$
is the final velocity.

\textbf{Deceleration} (retardation) is negative
acceleration --- the velocity is decreasing in magnitude.
\end{definitionbox}

\begin{misconceptionbox}
A body can be accelerating even when its speed is
constant. If a car goes around a circular bend at constant
speed, its direction is continuously changing, so its
velocity is changing, so it is accelerating. This
\textit{centripetal acceleration} is covered fully in
Chapter 6. For now, recognise that speed and velocity
are different, and acceleration means change in velocity
--- not necessarily change in speed.
\end{misconceptionbox}

% ============================================================
\section{Equations of Uniform Acceleration}
% ============================================================

When acceleration is \textbf{constant} (uniform), we can
derive four equations that connect displacement $s$,
initial velocity $u$, final velocity $v$, acceleration $a$,
and time $t$.

\begin{historybox}[Galileo and the Inclined Plane]
In the 1590s, Galileo Galilei rolled bronze balls down
smooth inclined planes in Pisa. He could not time free
fall directly --- it was too fast. By using a gentle slope,
he slowed the motion enough to time it with a water clock.
He discovered that the distance covered was proportional
to the square of the time --- $s \propto t^2$ --- the
mathematical signature of uniform acceleration. This was
the first experimental determination of the equations we
are about to derive, and it overturned 2000 years of
Aristotelian physics in one stroke.
\end{historybox}

\subsection{Derivation of the Four Equations}

\textbf{Equation 1: $v = u + at$}

From the definition of acceleration:
\[
  a = \frac{v - u}{t}
  \implies \boxed{v = u + at}
\]

\textbf{Equation 2: $s = \frac{1}{2}(u + v)t$}

For uniform acceleration, the average velocity is the
arithmetic mean of initial and final velocities:
\[
  \bar{v} = \frac{u + v}{2}
  \implies \boxed{s = \frac{u+v}{2}\cdot t}
\]

\textbf{Equation 3: $s = ut + \frac{1}{2}at^2$}

Substitute $v = u + at$ into Equation 2:
\[
  s = \frac{u + (u+at)}{2}\cdot t
  = \frac{2u + at}{2}\cdot t
  \implies \boxed{s = ut + \frac{1}{2}at^2}
\]

\textbf{Equation 4: $v^2 = u^2 + 2as$}

From Equation 1: $t = (v-u)/a$. Substitute into
Equation 2:
\[
  s = \frac{u+v}{2} \cdot \frac{v-u}{a}
  = \frac{v^2 - u^2}{2a}
  \implies \boxed{v^2 = u^2 + 2as}
\]

\begin{lawbox}[The Four Equations of Uniform Acceleration]
\begin{align}
  v &= u + at \label{eq:v}\\[4pt]
  s &= \frac{(u+v)}{2}\,t \label{eq:s_avg}\\[4pt]
  s &= ut + \tfrac{1}{2}at^2 \label{eq:s_quad}\\[4pt]
  v^2 &= u^2 + 2as \label{eq:v_sq}
\end{align}
Valid only when acceleration $a$ is \textbf{constant}.
Each equation contains four of the five variables
$\{u, v, a, s, t\}$. Choose the equation that contains
the three known variables and the one unknown.
\end{lawbox}

\vspace{3.3cm}

\begin{examtipbox}
Strategy for kinematics problems:
\begin{enumerate}[noitemsep]
  \item Write down what is given and what is unknown.
  \item Establish a sign convention (usually: positive
  direction = direction of initial motion).
  \item Choose the equation containing the unknown and
  the three knowns.
  \item Substitute carefully, paying attention to signs.
\end{enumerate}
\end{examtipbox}

\begin{examplebox}[3.1]
\stepcounter{workedexample}
A car starts from rest and accelerates uniformly at
$3.0\ \text{m\,s}^{-2}$ for $8.0\ \text{s}$.
Calculate: (a) the final velocity, (b) the distance
travelled.
\end{examplebox}

\begin{solutionbox}
Given: $u = 0$, $a = 3.0\ \text{m\,s}^{-2}$,
$t = 8.0\ \text{s}$\\
Unknown: $v$ and $s$

\textbf{(a)} Using $v = u + at$:
\[
  v = 0 + 3.0 \times 8.0 = 24\ \text{m\,s}^{-1}
\]

\textbf{(b)} Using $s = ut + \frac{1}{2}at^2$:
\[
  s = 0 \times 8.0 + \tfrac{1}{2} \times 3.0 \times 8.0^2
  = 0 + \tfrac{1}{2} \times 3.0 \times 64
  = 96\ \text{m}
\]

\textbf{Check} using $v^2 = u^2 + 2as$:
$24^2 = 0 + 2(3.0)(96) = 576$ ✓
\end{solutionbox}

\begin{examplebox}[3.2]
\stepcounter{workedexample}
A lorry travelling at $25\ \text{m\,s}^{-1}$ applies its
brakes and decelerates uniformly. It comes to rest after
travelling $125\ \text{m}$. Find: (a) the deceleration,
(b) the time taken to stop.
\end{examplebox}

\begin{solutionbox}
Given: $u = 25\ \text{m\,s}^{-1}$,
$v = 0$, $s = 125\ \text{m}$

\textbf{(a)} Using $v^2 = u^2 + 2as$:
\[
  0 = 25^2 + 2a(125)
  \implies 250a = -625
  \implies a = -2.5\ \text{m\,s}^{-2}
\]
Deceleration $= 2.5\ \text{m\,s}^{-2}$

\textbf{(b)} Using $v = u + at$:
\[
  0 = 25 + (-2.5)t
  \implies t = \frac{25}{2.5} = 10\ \text{s}
\]

\begin{realworldbox}[Braking Distance on Nigerian Roads]
The lorry in this example requires $125\ \text{m}$ to
stop from $25\ \text{m\,s}^{-1}$ ($90\ \text{km\,h}^{-1}$)
--- that is over a third of a football pitch. On wet or
laterite roads, the deceleration is smaller, and the
stopping distance is even longer. This is why heavy goods
vehicles are deadly in urban traffic: their stopping
distances are enormous.
\end{realworldbox}
\end{solutionbox}

% ============================================================
\section{Motion Graphs}
% ============================================================

Graphs are among the most powerful tools for analysing
motion. Two graphs are used routinely: the
\textbf{displacement-time} graph and the
\textbf{velocity-time} graph.

\subsection{Displacement-Time Graphs}

\begin{lawbox}[Displacement-Time Graph]
\begin{itemize}[noitemsep]
  \item The \textbf{gradient} (slope) of a
  displacement-time graph equals the \textbf{velocity}
  at that instant.
  \item A straight line $\implies$ constant velocity
  (uniform motion).
  \item A curved line $\implies$ changing velocity
  (acceleration or deceleration).
  \item Horizontal line $\implies$ object at rest ($v = 0$).
  \item Negative gradient $\implies$ body moving in the
  negative direction.
\end{itemize}
\end{lawbox}

\begin{diagbox}[Displacement-Time Graphs]
\begin{center}
\begin{tikzpicture}
\begin{axis}[
  width=12cm, height=5cm,
  xlabel={Time $t$ (s)},
  ylabel={Displacement $s$ (m)},
  xmin=0, xmax=6, ymin=0, ymax=30,
  xtick={0,1,2,3,4,5,6},
  ytick={0,5,10,15,20,25,30},
  grid=major, grid style={gray!30},
 legend style={
  at={(0.5,1.12)},
  anchor=south,
  legend columns=3,
  font=\small
},
]
  % Constant velocity (straight line)
  \addplot[codexblue, thick, domain=0:4]
    {5*x} node[right, pos=0.7]{};
  \addlegendentry{Constant velocity}

  % Acceleration (curve)
  \addplot[codexred, thick, domain=0:5, samples=50]
    {x^2} node[right, pos=0.8]{};
  \addlegendentry{Acceleration (slope increasing)}

  % At rest (horizontal)
  \addplot[codexgreen, thick, domain=0:6]
    {20} node[right, pos=0.9]{};
  \addlegendentry{At rest}
\end{axis}
\end{tikzpicture}
\end{center}
\end{diagbox}

\subsection{Velocity-Time Graphs}

\begin{lawbox}[Velocity-Time Graph]
\begin{itemize}[noitemsep]
  \item The \textbf{gradient} of a velocity-time graph
  equals the \textbf{acceleration} at that instant.
  \item The \textbf{area under} a velocity-time graph
  (between the graph and the time axis) equals the
  \textbf{displacement}.
  \item Straight line with positive gradient
  $\implies$ uniform acceleration.
  \item Straight line with negative gradient
  $\implies$ uniform deceleration.
  \item Horizontal line $\implies$ constant velocity
  ($a = 0$).
\end{itemize}
\end{lawbox}

\begin{diagbox}[Velocity-Time Graph]
\begin{center}
\begin{tikzpicture}
\begin{axis}[
  width=12cm, height=6cm,
  xlabel={Time $t$ (s)},
  ylabel={Velocity $v$ (m\,s$^{-1}$)},
  xmin=0, xmax=12, ymin=0, ymax=30,
  xtick={0,2,4,6,8,10,12},
  ytick={0,5,10,15,20,25,30},
  grid=major, grid style={gray!30},
]
  % Phase 1: Acceleration 0 to 4s, 0 to 20 m/s
  \addplot[codexgreen, very thick, domain=0:4]{5*x};
  \node[codexgreen, font=\small] at (axis cs:2,13)
    {Accelerating};

  % Phase 2: Constant velocity 4s to 8s
  \addplot[codexblue, very thick, domain=4:8]{20};
  \node[codexblue, font=\small] at (axis cs:6,22)
    {Constant velocity};

  % Phase 3: Deceleration 8s to 12s
  \addplot[codexred, very thick, domain=8:12]{20-5*(x-8)};
  \node[codexred, font=\small] at (axis cs:10,13)
    {Decelerating};

  % Area shading
  \addplot[fill=codexgreen!15, opacity=0.5, domain=0:4]
    {5*x} \closedcycle;
  \addplot[fill=codexblue!15, opacity=0.5, domain=4:8]
    {20} \closedcycle;
  \addplot[fill=codexred!15, opacity=0.5, domain=8:12]
    {20-5*(x-8)} \closedcycle;

  % Rise/run triangle: the slope of a v-t graph is acceleration.
  \draw[dashed, gray, thick]
    (axis cs:1,5) -- (axis cs:2,5) -- (axis cs:2,10);
  \node[below, gray, font=\scriptsize] at (axis cs:1.5,5) {$\Delta t$};
  \node[right, gray, font=\scriptsize] at (axis cs:2,7.5) {$\Delta v$};
\end{axis}
\end{tikzpicture}
\end{center}
\small\textit{The shaded areas represent displacement in
each phase.}
\end{diagbox}

\begin{examplebox}[3.3]
\stepcounter{workedexample}
From the velocity-time graph above, calculate:
(a) the acceleration in phase 1;
(b) the total displacement over all three phases;
(c) the deceleration in phase 3.
\end{examplebox}

\begin{solutionbox}
\textbf{(a)} Phase 1 gradient (0 to 4 s):
\[
  a = \frac{\Delta v}{\Delta t}
  = \frac{20 - 0}{4 - 0} = 5\ \text{m\,s}^{-2}
\]

\textbf{(b)} Total displacement = total area under graph:
\begin{align*}
  s_1 &= \tfrac{1}{2} \times 4 \times 20 = 40\ \text{m}
  \quad\text{(triangle)}\\
  s_2 &= 4 \times 20 = 80\ \text{m}
  \quad\text{(rectangle)}\\
  s_3 &= \tfrac{1}{2} \times 4 \times 20 = 40\ \text{m}
  \quad\text{(triangle)}\\
  s_{total} &= 40 + 80 + 40 = 160\ \text{m}
\end{align*}

\textbf{(c)} Phase 3 gradient (8 to 12 s):
\[
  a = \frac{0 - 20}{12 - 8} = -5\ \text{m\,s}^{-2}
\]
Deceleration $= 5\ \text{m\,s}^{-2}$
\end{solutionbox}

\vspace{3.3cm}

% ============================================================
\section{Free Fall Under Gravity}
% ============================================================

\begin{definitionbox}[Free Fall]
\textbf{Free fall} is motion under the influence of
gravity alone, with no air resistance. Near the surface
of the Earth, every freely falling object accelerates
downward at:
\[
  g = 9.81\ \text{m\,s}^{-2}
  \approx 10\ \text{m\,s}^{-2}\ \text{(UTME)}
\]
This value is the same for all objects, regardless of
their mass. A feather and a hammer, dropped together in
a vacuum, fall at exactly the same rate.
\end{definitionbox}


\begin{historybox}[Galileo Drops From the Tower]
Legend says Galileo dropped two cannonballs of different
masses from the Leaning Tower of Pisa to disprove
Aristotle's claim that heavier objects fall faster.
Whether or not the story is literally true, Galileo
certainly established by careful experiment that all
objects fall with the same acceleration in the absence
of air resistance. This was revolutionary: for 2000
years, educated people had believed otherwise.
\end{historybox}

For vertical motion under gravity, the equations of
uniform acceleration apply with $a = g$ (taking downward
as positive) or $a = -g$ (taking upward as positive).
We will use the convention: \textbf{upward positive,
downward negative.}

\begin{align*}
  v &= u - gt \quad \text{(upward positive)}\\
  h &= ut - \tfrac{1}{2}gt^2\\
  v^2 &= u^2 - 2gh
\end{align*}

\begin{examplebox}[3.4]
\stepcounter{workedexample}
A stone is thrown vertically upward from the ground
with an initial velocity of $20\ \text{m\,s}^{-1}$.
Taking $g = 10\ \text{m\,s}^{-2}$, find:
(a) the maximum height reached;
(b) the time taken to reach maximum height;
(c) the total time before the stone returns to the
ground;
(d) the velocity when it returns to ground level.
\end{examplebox}

\begin{solutionbox}
Taking upward as positive: $u = +20\ \text{m\,s}^{-1}$,
$a = -10\ \text{m\,s}^{-2}$.

\textbf{(a)} At maximum height, $v = 0$. Using
$v^2 = u^2 + 2as$:
\[
  0 = 20^2 + 2(-10)s
  \implies 20s = 400
  \implies s = H_{max} = 20\ \text{m}
\]

\textbf{(b)} Time to maximum height. Using $v = u + at$:
\[
  0 = 20 + (-10)t \implies t = 2\ \text{s}
\]

\textbf{(c)} By symmetry, the time to fall back from
maximum height to the ground equals the time to rise:
\[
  t_{total} = 2 \times 2 = 4\ \text{s}
\]

\textbf{(d)} Using $v = u + at$ for full journey:
\[
  v = 20 + (-10)(4) = 20 - 40 = -20\ \text{m\,s}^{-1}
\]
The negative sign confirms the velocity is downward.
The speed of impact equals the launch speed --- the
motion is perfectly symmetric (no air resistance).
\end{solutionbox}

\begin{examplebox}[3.5]
\stepcounter{workedexample}
A coconut falls from a palm tree $20\ \text{m}$ tall.
(a) How long does it take to reach the ground?
(b) With what velocity does it strike the ground?
($g = 10\ \text{m\,s}^{-2}$, neglect air resistance)
\end{examplebox}

\begin{solutionbox}
Taking downward as positive: $u = 0$, $a = +10\ \text{m\,s}^{-2}$,
$s = 20\ \text{m}$.

\textbf{(a)} Using $s = ut + \frac{1}{2}at^2$:
\[
  20 = 0 + \tfrac{1}{2}(10)t^2
  = 5t^2
  \implies t^2 = 4
  \implies t = 2\ \text{s}
\]

\textbf{(b)} Using $v = u + at$:
\[
  v = 0 + 10 \times 2 = 20\ \text{m\,s}^{-1}
  \text{ (downward)}
\]

\begin{realworldbox}[Coconut Danger]
Coconut trees can be over 25 m tall. A coconut falling
from 25 m strikes the ground at approximately
$22\ \text{m\,s}^{-1}$ ($79\ \text{km\,h}^{-1}$).
This is why the area beneath tall coconut trees is
genuinely dangerous, and why plantation workers in
parts of Nigeria and across coastal West Africa use
protective gear when harvesting.
\end{realworldbox}
\end{solutionbox}

% ============================================================
\section{Projectile Motion}
% ============================================================

A \textbf{projectile} is any object that is given an
initial velocity and then moves freely under gravity
alone. Thrown stones, footballs, artillery shells, and
water from a hose are all projectiles.

The key insight for analysing projectile motion is this:
\textbf{the horizontal and vertical motions are completely
independent of each other.}

\begin{keybox}[Independence of Horizontal and Vertical Motion]
\begin{itemize}[noitemsep]
  \item \textbf{Horizontal:} No force acts horizontally
  (ignoring air resistance). Horizontal velocity is
  \textit{constant} throughout the flight.
  \[
    v_x = u\cos\theta = \text{constant}
    \qquad x = u\cos\theta \cdot t
  \]
  \item \textbf{Vertical:} Gravity acts downward.
  The vertical motion is free fall.
  \[
    v_y = u\sin\theta - gt
    \qquad y = u\sin\theta \cdot t - \tfrac{1}{2}gt^2
  \]
\end{itemize}
\end{keybox}

\begin{diagbox}[Projectile Motion Trajectory]
\begin{center}
\begin{tikzpicture}
\begin{axis}[
  width=12cm, height=7cm,
  xlabel={Horizontal distance $x$ (m)},
  ylabel={Height $y$ (m)},
  xmin=0, xmax=50, ymin=0, ymax=15,
  grid=major, grid style={gray!20},
  axis equal=false,
  clip=false,
]
  % Trajectory: u=20, theta=45, g=10
  % x = 20cos45*t, y = 20sin45*t - 5t^2
  % T = 2*20sin45/10 = 2.83 s, R = 40 m
  \addplot[codexblue, very thick, domain=0:2.828, samples=60,
    variable=\t]
    ({20*0.7071*\t}, {20*0.7071*\t - 5*\t*\t});
  \addlegendentry{$u=20$ m/s, $\theta=45°$}

  % Sample positions along the trajectory (not velocity vectors).
  \foreach \tt/\vx/\vy in {
    0.5/{7.07}/{7.07-5},
    1.0/{7.07}/{7.07-10},
    1.5/{7.07}/{7.07-15},
    2.0/{7.07}/{7.07-20}
  }{
    \pgfmathsetmacro\px{20*0.7071*\tt}
    \pgfmathsetmacro\py{20*0.7071*\tt - 5*\tt*\tt}
    \edef\temp{\noexpand\fill[codexred] (axis cs:\px,\py)
      circle (1.5pt);}
    \temp
  }
  \addlegendimage{only marks, mark=*, mark options={draw=codexred, fill=codexred}}
  \addlegendentry{Sample positions}

  % Horizontal component arrow
  \draw[->, codexred, thick]
    (axis cs:20,11.2) -- (axis cs:23,11.2)
    node[right, font=\tiny]{$v_x = \text{const}$};

  % Max height annotation
  \draw[dashed, gray] (axis cs:20,0) -- (axis cs:20,10);
  \node[right, gray, font=\small] at (axis cs:20,5)
    {$H_{max}$};

  % Range annotation
  \draw[<->, gray] (axis cs:0,-3.2) -- (axis cs:40,-3.2);
  \node[below, gray, font=\small] at (axis cs:20,-3.2)
    {Range $R$};
\end{axis}
\end{tikzpicture}
\end{center}
\end{diagbox}

\subsection{Key Projectile Formulae}

For a projectile launched at speed $u$ at angle $\theta$
above the horizontal:

\begin{lawbox}[Projectile Equations (launch from ground)]
\begin{align*}
  t &= \frac{u\sin\theta}{g}
  \quad \text{(time taken to reach the maximum height)}\\[6pt]
  T &= \frac{2u\sin\theta}{g}
  \quad \text{(time of flight)}\\[6pt]
  H_{max} &= \frac{u^2\sin^2\theta}{2g}
  \quad \text{(maximum height)}\\[6pt]
  R &= \frac{u^2\sin 2\theta}{g}
  \quad \text{(horizontal range)}\\[6pt]
  v_x &= u\cos\theta
  \quad \text{(horizontal velocity, constant)}\\[6pt]
  v_y &= u\sin\theta - gt
  \quad \text{(vertical velocity at time } t\text{)}
\end{align*}
Maximum range occurs at $\theta = 45°$.
\end{lawbox}

\begin{lawbox}[Horizontal Projection from Height $H$]
If a projectile is launched \textit{horizontally} (i.e.\
$\theta = 0$) from a height $H$:
\[
  t_{fall} = \sqrt{\frac{2H}{g}}
  \quad \text{(time to hit ground)}
  \qquad
  R = u\sqrt{\frac{2H}{g}}
  \quad \text{(range)}
\]
\end{lawbox}

\begin{examplebox}[3.6]
\stepcounter{workedexample}
A footballer kicks a ball at $18\ \text{m\,s}^{-1}$ at
an angle of $40°$ above the ground. Taking
$g = 10\ \text{m\,s}^{-2}$, find:
(a) the maximum height the ball reaches;
(b) the total time the ball is in the air;
(c) the horizontal range.
\end{examplebox}

\begin{solutionbox}
Given: $u = 18\ \text{m\,s}^{-1}$, $\theta = 40°$,
$g = 10\ \text{m\,s}^{-2}$.

\[
  u_x = 18\cos 40° = 18 \times 0.766 = 13.8\ \text{m\,s}^{-1}
\]
\[
  u_y = 18\sin 40° = 18 \times 0.643 = 11.6\ \text{m\,s}^{-1}
\]

\textbf{(a)} Maximum height:
\[
  H_{max} = \frac{u_y^2}{2g}
  = \frac{11.6^2}{2 \times 10}
  = \frac{134.56}{20} = 6.7\ \text{m}
\]

\textbf{(b)} Time of flight:
\[
  T = \frac{2u_y}{g}
  = \frac{2 \times 11.6}{10} = 2.32\ \text{s}
\]

\textbf{(c)} Horizontal range:
\[
  R = u_x \times T = 13.8 \times 2.32 = 32.0\ \text{m}
\]
\end{solutionbox}

\begin{examplebox}[3.7]
\stepcounter{workedexample}
A stone is thrown horizontally at $15\ \text{m\,s}^{-1}$
from the top of a cliff $45\ \text{m}$ high.
Find: (a) the time it takes to reach the sea;
(b) the horizontal distance from the base of the cliff.
\end{examplebox}

\begin{solutionbox}
Horizontal: $u_x = 15\ \text{m\,s}^{-1}$, no acceleration.\\
Vertical: $u_y = 0$, $a = 10\ \text{m\,s}^{-2}$ downward,
$H = 45\ \text{m}$.

\textbf{(a)} Time to fall $45\ \text{m}$:
\[
  H = \tfrac{1}{2}gt^2
  \implies 45 = \tfrac{1}{2}(10)t^2
  = 5t^2
  \implies t = \sqrt{9} = 3\ \text{s}
\]

\textbf{(b)} Horizontal distance:
\[
  x = u_x \times t = 15 \times 3 = 45\ \text{m}
\]
\end{solutionbox}

\vspace{3.3cm}

% ============================================================
\section{Chapter Summary}
% ============================================================

\begin{summarybox}

\textbf{Distance} (scalar) vs \textbf{displacement}
(vector). \textbf{Speed} (scalar) vs \textbf{velocity}
(vector).

\vspace{0.3cm}

\textbf{Four equations of uniform acceleration:}
\begin{align*}
  v &= u + at\\
  s &= \tfrac{1}{2}(u+v)t\\
  s &= ut + \tfrac{1}{2}at^2\\
  v^2 &= u^2 + 2as
\end{align*}

\textbf{Velocity-time graph:}
gradient = acceleration; area = displacement.

\textbf{Free fall:} $a = g \approx 10\ \text{m\,s}^{-2}$
downward for all objects, mass independent.

\textbf{Projectile motion:} horizontal and vertical
motions are independent. Horizontal velocity is constant.
\[
  T = \frac{2u\sin\theta}{g}, \quad
  H = \frac{u^2\sin^2\theta}{2g}, \quad
  R = \frac{u^2\sin 2\theta}{g}
\]
Maximum range at $\theta = 45°$.

\end{summarybox}

% ============================================================
% EXERCISES
% ============================================================
\newpage
\begin{exercisebox}[3]

\subsection*{Section A: Multiple Choice}

\textbf{A1.} A body starts from rest and accelerates
uniformly. After 5 s it has travelled 100 m. Its
acceleration is:

\mcoption{A}{$4\ \text{m\,s}^{-2}$}
\mcoption{B}{$8\ \text{m\,s}^{-2}$}
\mcoption{C}{$10\ \text{m\,s}^{-2}$}
\mcoption{D}{$20\ \text{m\,s}^{-2}$}
\ansref{Ch3-A1}

\vspace{0.4cm}

\textbf{A2.} On a velocity-time graph, a horizontal
straight line indicates:

\mcoption{A}{Zero velocity}
\mcoption{B}{Constant acceleration}
\mcoption{C}{Constant velocity}
\mcoption{D}{Increasing acceleration}
\ansref{Ch3-A2}

\vspace{0.4cm}

\textbf{A3.} A ball is thrown vertically upward with
speed $u$. The ratio of its speed at half the maximum
height to its speed at the maximum height is:

\mcoption{A}{$\sqrt{2} : 1$}
\mcoption{B}{$1 : 0$}
\mcoption{C}{$\sqrt{2} : 0$}
\mcoption{D}{$1 : \sqrt{2}$}
\ansref{Ch3-A3}

\vspace{0.4cm}

\textbf{A4.} A projectile is fired at $45°$ to the
horizontal with speed $v$. Its range is $R$. If the
same projectile is fired with the same speed at $30°$,
its new range is approximately:

\mcoption{A}{$0.5R$}
\mcoption{B}{$0.87R$}
\mcoption{C}{$R$}
\mcoption{D}{$1.15R$}
\ansref{Ch3-A4}

\vspace{0.4cm}

\textbf{A5.} A stone is dropped from a height $H$ and
takes time $T$ to reach the ground. A second stone is
dropped from height $4H$. The time it takes is:

\mcoption{A}{$T$}
\mcoption{B}{$2T$}
\mcoption{C}{$4T$}
\mcoption{D}{$16T$}
\ansref{Ch3-A5}

\vspace{0.4cm}

\textbf{A6.} The area under a velocity-time graph
between $t = 0$ and $t = T$ represents:

\mcoption{A}{Acceleration}
\mcoption{B}{Distance only}
\mcoption{C}{Displacement}
\mcoption{D}{The product of average velocity and
acceleration}
\ansref{Ch3-A6}

\vspace{0.4cm}

\textbf{A7.} A ball is thrown horizontally from a tower.
Ignoring air resistance, which statement is correct?

\mcoption{A}{Its horizontal speed decreases as it falls}
\mcoption{B}{Its vertical speed is constant}
\mcoption{C}{Its horizontal and vertical motions are
independent}
\mcoption{D}{It reaches maximum height at the midpoint
of its flight}
\ansref{Ch3-A7}

\vspace{0.4cm}

\textbf{A8.} A car travels $60\ \text{m}$ in the first
$4\ \text{s}$ of its journey and $100\ \text{m}$ in the
next $4\ \text{s}$. The car is:

\mcoption{A}{Decelerating}
\mcoption{B}{Moving at constant velocity}
\mcoption{C}{Accelerating}
\mcoption{D}{Moving in a circle}
\ansref{Ch3-A8}

\vspace{0.4cm}

\textbf{A9.} A body moving at $30\ \text{m\,s}^{-1}$
decelerates uniformly and stops after $6\ \text{s}$.
The distance it travels during deceleration is:

\mcoption{A}{$45\ \text{m}$}
\mcoption{B}{$90\ \text{m}$}
\mcoption{C}{$135\ \text{m}$}
\mcoption{D}{$180\ \text{m}$}
\ansref{Ch3-A9}

\vspace{0.4cm}

\textbf{A10.} Which pair of physical quantities has the
same unit?

\mcoption{A}{Speed and acceleration}
\mcoption{B}{Displacement and distance}
\mcoption{C}{Velocity and speed}
\mcoption{D}{Distance and time}
\ansref{Ch3-A10}

\vspace{0.6cm}

\subsection*{Section B: Theory and Calculations}

\textbf{B1.} A train starts from rest at a station and
accelerates uniformly at $0.8\ \text{m\,s}^{-2}$ for
$30\ \text{s}$. It then travels at constant velocity for
$2\ \text{min}$. Finally, it decelerates uniformly and
comes to rest at the next station in $20\ \text{s}$.\\[4pt]
(a) Calculate the maximum velocity reached.\\
(b) Sketch the velocity-time graph for the entire journey,
labelling all values.\\
(c) Calculate the total distance between the two
stations.\\
(d) Calculate the deceleration in the final phase.
\hfill\ansref{Ch3-B1}

\vspace{0.4cm}

\textbf{B2.} A rock is dropped from a bridge into a
river below. The splash is heard $2.8\ \text{s}$ after
the rock is released. The speed of sound in air is
$340\ \text{m\,s}^{-1}$ and $g = 10\ \text{m\,s}^{-2}$.\\[4pt]
(a) Write an equation relating the height $H$ of the
bridge to the fall time $t_1$ and the sound travel
time $t_2$, given that $t_1 + t_2 = 2.8\ \text{s}$.\\
(b) Solve for $H$ (this requires solving a quadratic
equation).\\
(c) With what velocity does the rock hit the water?
\hfill\ansref{Ch3-B2}

\vspace{0.4cm}

\textbf{B3.} A goalkeeper kicks a football from ground
level at a speed of $22\ \text{m\,s}^{-1}$ at an angle
of $35°$ above the horizontal. ($g = 10\ \text{m\,s}^{-2}$)\\[4pt]
(a) Find the horizontal and vertical components of the
initial velocity.\\
(b) Calculate the maximum height reached.\\
(c) Calculate the total time the ball is in the air.\\
(d) Calculate the horizontal range.\\
(e) A defender stands $35\ \text{m}$ away from the
goalkeeper. At what height is the ball when it passes
over the defender?
\hfill\ansref{Ch3-B3}

\vspace{0.4cm}

\textbf{B4.} From a cliff $80\ \text{m}$ above sea level,
a rescue flare is fired \textit{vertically upward} at
$30\ \text{m\,s}^{-1}$. Taking $g = 10\ \text{m\,s}^{-2}$
and setting the base of the cliff as $y = 0$:
(a) Write an equation for the height $y$ of the flare
above sea level as a function of time.\\
(b) Find the maximum height above sea level.\\
(c) Find the time at which the flare hits the sea.\\
(d) Find the speed at which the flare hits the sea.
\hfill\ansref{Ch3-B4}

\vspace{0.4cm}

\textbf{B5.} Two motorcyclists, Adebisi and Chukwudi,
start from the same point. Adebisi rides due north at
constant $60\ \text{km\,h}^{-1}$. At the same instant,
Chukwudi rides due east at constant $80\ \text{km\,h}^{-1}$.\\[4pt]
(a) After $30\ \text{min}$, how far apart are they?\\
(b) What is the velocity of Adebisi relative to Chukwudi
(magnitude and direction)?\\
(c) At what rate is the distance between them increasing
after $30\ \text{min}$?
\hfill\ansref{Ch3-B5}

\end{exercisebox}
